\documentclass[10pt,a4paper]{amsart}
\usepackage[margin=19mm]{geometry}
\usepackage{amsmath,amssymb,mathtools}
\usepackage[scr=rsfs,cal=euler]{mathalfa}
\usepackage{xfrac}
\usepackage[unicode,hidelinks,bookmarks=false]{hyperref}
\newcommand{\ie}{i.e.\ }
\newcommand{\cf}{cf.\ }
\newcommand{\resp}{respectively\ }
\newcommand{\Subsection}{\subsection}
\providecommand{\nobreakdash}{\nobreak\hbox{-}}
\setlength{\emergencystretch}{2em}
\multlinegap0pt
\usepackage[shortlabels]{enumitem}
\usepackage{amssymb}
\usepackage{mathtools}
\usepackage{xcolor}
\usepackage{dsfont}
\usepackage{float}
\newtheorem{lemma}{Lemma}
\newtheorem{definition}{Definition}
\newtheorem{proposition}{Proposition}
\usepackage{cleveref}
\crefname{lemma}{Lemma}{Lemmas}
\crefname{definition}{Definition}{Definitions}
\crefname{proposition}{Proposition}{Propositions}
\newcommand{\calP}{\mathcal P}
\title{Nonparametric sequential change detection without partitioning}
\author{Aytijhya Saha, Aaditya Ramdas}
\date{Source: arXiv:2607.28322v1 (CC BY 4.0)}
\begin{document}
\maketitle

\noindent\emph{Source and licence.} Nonparametric sequential change detection without partitioning. Version arXiv:2607.28322v1 (CC BY 4.0), distributed by arXiv under the Creative Commons Attribution 4.0 International licence, http://creativecommons.org/licenses/by/4.0/, as recorded in the arXiv licence metadata for this exact version and shown on the abstract page https://arxiv.org/abs/2607.28322v1. Full licence notice: \texttt{licenses/arxiv-2607.28322-CC-BY-4.0.txt}.\par\smallskip

\noindent\textit{Editorial scope.} Counted unit: the article's proposition prop:gaussian-local-witness (source0.tex lines 1293--1296). The article's complete original proof of it is delivered with it (source0.tex lines 5539--6236, one \texttt{proof} environment of 698 lines, which the article places away from the statement). No mathematics is written, repaired or re-derived: the statement and the proof are the article's own lines, in the article's own notation, and no retained line is substituted or re-flowed.  Retained uncounted as the source-local context the proof needs: lines 1136--1142; lines 1240--1247; lines 1253--1276.  Nothing was omitted from any retained span, so the package carries no omission record.  No retained line contains a citation, so the wrapper carries no bibliography and no bibliography entry.  No retained line contains a figure environment, an image-inclusion command or drawing code, so no image is delivered and the package declares no asset dependency.  Editorial formatting only, in addition: an AMS \texttt{amsart} preamble that restates the article's own declarations and macros used by the retained lines, namely the environments \texttt{lemma}, \texttt{definition}, \texttt{proposition} on separate counters, together with the standard packages those lines need.  The retained environments print in this document as Lemma 1, Definition 1, Proposition 1 in order of appearance, whereas the article prints the same items with the numbering of its own full document.  The complete native source of the article as distributed in the arXiv e-print for this exact version is bundled unchanged as \texttt{sources/206407/source0.tex}.
\begin{lemma}[Adjusted running maximum is an e-process]\label{lem:adjusted-max}
Let $(E_t)_{t\ge0}$ be an e-process under a null law $\mathbb P$, with $E_0=1$, and let
\[
  E_t^*=\max_{0\le u\le t}E_u.
\]
If $a$ is an adjuster, then $(a(E_t^*))_{t\ge0}$ is an e-process under $\mathbb P$.
\end{lemma}

\begin{lemma}[Weakly compact classes are simultaneously REGROW-regular]
\label{lem:weakly-compact-regrow}
Let $\mathcal X$ be a Polish space, and let
$ K\subseteq\mathcal M_1(\mathcal X)$
be a nonempty weakly compact class of probability measures.
Then $K$ is {simultaneously REGROW-regular over any
$\mathcal A \subseteq K^c$}. 
\end{lemma}

\begin{definition}[Countable local REGROW witness basis]\label{def:witness-basis}
We say that $\calP$ has a countable local REGROW witness basis if there exists a countable family $\mathscr B=\{B_j:j\ge 1\}$ of relatively weakly open subsets of $\calP$ such that:
\begin{enumerate}[label=\textup{(\roman*)}]
\item
For every $P\in\calP$ and every relatively weakly open
neighborhood $G$ of $P$, there exists $j\ge1$ such that
\[
    P\in B_j
    \subseteq
    \overline B_j
    \subseteq
    G,
\]
where $\overline B_j$ denotes relative weak closure in $\calP$.

\item
For every $j\ge1$, the class $\overline B_j$ is weakly compact. 
\item
For every $j\ge1$, the exterior class
$B_j^c:=\calP\setminus B_j$
is simultaneously REGROW-regular over $B_j$. 
\end{enumerate}
Moreover, $\mathscr B$ is called the countable local REGROW witness basis of $\calP$.
\end{definition}

\begin{proposition}[Gaussian location family and REGROW witness]
\label{prop:gaussian-local-witness}
Let $\mathcal{P} = \{N(\theta, 1) : \theta \in \mathbb{R}\}$ be the unit-variance Gaussian location family. Then $\mathcal{P}$ is not weakly compact, but it admits a countable local REGROW witness basis in the sense of Definition \ref{def:witness-basis}.
\end{proposition}

\begin{proof}[Proof of \Cref{prop:gaussian-local-witness}]
We divide the proof into two parts.
\paragraph{Part 1: $\mathcal P$ is not weakly compact.}

By Prokhorov's theorem, a family of probability measures on a
Polish space is relatively weakly compact if and only if it is
uniformly tight.

Let $K\subseteq\mathbb R$ be compact. Then there exists $M>0$
such that
\[
    K\subseteq[-M,M].
\]
Consider the sequence
\[
    P_n=N(n,1)\in\mathcal P,
    \qquad n\ge1.
\]
Writing $\Phi$ for the standard Gaussian distribution function,
we have
\[
\begin{aligned}
    P_n(K)
    &\le
    P_n([-M,M])\\
    &=
    \Phi(M-n)-\Phi(-M-n)
    \longrightarrow0
\end{aligned}
\]
as $n\to\infty$. Consequently, for every compact
$K\subseteq\mathbb R$,
\[
    \inf_{P\in\mathcal P}P(K)=0.
\]
In particular, there is no compact $K$ such that
\[
    \inf_{P\in\mathcal P}P(K)\ge\frac12.
\]
Thus $\mathcal P$ is not uniformly tight. By Prokhorov's
theorem, $\mathcal P$ is not relatively weakly compact and hence
is not weakly compact.

\paragraph{Part 2: Construction of a countable local
REGROW witness basis.}

Consider the parameterization
\[
    f:\mathbb R\longrightarrow\mathcal P,
    \qquad
    f(\theta)=P_\theta.
\]
We first verify that $f$ is a homeomorphism from $\mathbb R$,
equipped with its usual topology, onto $\mathcal P$, equipped
with the relative weak topology.

If $\theta_n\to\theta$, then the characteristic functions satisfy
\[
    \exp\left\{
        i u\theta_n-\frac{u^2}{2}
    \right\}
    \longrightarrow
    \exp\left\{
        i u\theta-\frac{u^2}{2}
    \right\}
\]
for every $u\in\mathbb R$. Hence, by L\'evy's continuity theorem,
\[
    P_{\theta_n}\Rightarrow P_\theta.
\]
Thus $f$ is continuous.

Conversely, suppose
\[
    P_{\theta_n}\Rightarrow P_\theta.
\]
Because the distribution function of $P_\theta$ is continuous
at $\theta$, weak convergence implies
\[
\begin{aligned}
    \Phi(\theta-\theta_n)
    &=
    P_{\theta_n}((-\infty,\theta])\\
    &\longrightarrow
    P_\theta((-\infty,\theta])
    =
    \frac12.
\end{aligned}
\]
Since $\Phi$ is continuous and strictly increasing,
\[
    \theta-\theta_n\longrightarrow0,
\]
and therefore $\theta_n\to\theta$. Since the weak topology on
$\mathcal P$ is metrizable, this proves continuity of $f^{-1}$.
Hence $f$ is a homeomorphism.

Define
\[
    \mathscr B
    :=
    \left\{
        B_{q,r}:
        q,r\in\mathbb Q,\ q<r
    \right\},
\]
where
\[
    B_{q,r}
    :=
    \{P_\theta\in\mathcal P:q<\theta<r\}.
\]
The collection $\mathscr B$ is countable. We verify the three
conditions in \Cref{def:witness-basis}.

\medskip

\noindent
\emph{Condition \textup{(i)}: local refinement.}

Fix $P_{\theta_0}\in\mathcal P$, and let $G$ be a relatively
weakly open neighborhood of $P_{\theta_0}$. Since $f$ is a
homeomorphism,
\[
    f^{-1}(G)
\]
is an open neighborhood of $\theta_0$ in $\mathbb R$. Hence
there exists $\epsilon>0$ such that
\[
    [\theta_0-\epsilon,\theta_0+\epsilon]
    \subseteq
    f^{-1}(G).
\]
By density of the rational numbers, we can choose
$q,r\in\mathbb Q$ such that
\[
    \theta_0-\epsilon<q<\theta_0<r<\theta_0+\epsilon.
\]
Then
\[
    P_{\theta_0}\in B_{q,r}.
\]
Moreover, because $f$ is a homeomorphism,
\[
    \overline{B_{q,r}}^{\,\mathcal P}
    =
    \{P_\theta:q\le\theta\le r\}
    =
    f([q,r]),
\]
where the closure is taken relative to $\mathcal P$. Therefore,
\[
    P_{\theta_0}
    \in
    B_{q,r}
    \subseteq
    \overline{B_{q,r}}^{\,\mathcal P}
    \subseteq
    G.
\]
Thus Condition~\textup{(i)} holds.

\medskip

\noindent
\emph{Condition \textup{(ii)}: simultaneous inner witnesses.}

Fix $q,r\in\mathbb Q$ with $q<r$. The relative weak closure of
$B_{q,r}$ is
\[
    K_{q,r}^{\mathrm{in}}
    :=
    \overline{B_{q,r}}^{\,\mathcal P}
    =
    \{P_\theta:q\le\theta\le r\}
    =
    f([q,r]).
\]
Since $[q,r]$ is compact and $f$ is continuous,
$K_{q,r}^{\mathrm{in}}$ is weakly compact as a subset of
$\mathcal M_1(\mathbb R)$.

By \Cref{lem:weakly-compact-regrow}, there exists a single
e-process
\[
    V^{q,r,\mathrm{in}}
    =
    \{V_t^{q,r,\mathrm{in}}\}_{t\ge0}
\]
which is valid under every $R\in K_{q,r}^{\mathrm{in}}$ and
satisfies, for every
$Q\notin K_{q,r}^{\mathrm{in}}$,
\[
    \liminf_{t\to\infty}
    \frac1t
    \log V_t^{q,r,\mathrm{in}}
    \ge
    \Phi_{K_{q,r}^{\mathrm{in}}}(Q),
    \qquad
    Q^\infty\text{-a.s.}
\]
Importantly, the same process $V^{q,r,\mathrm{in}}$ works for
every $Q\notin K_{q,r}^{\mathrm{in}}$.

If $V^{q,r,\mathrm{in}}$ is not already nondecreasing, let $a$
be a fixed growth-preserving e-process adjuster and define
\[
    \bigl(V^{q,r,\mathrm{in}}\bigr)_t^*
    :=
    \max_{0\le u\le t}
    V_u^{q,r,\mathrm{in}},
\]
and
\[
    E_t^{q,r,\mathrm{in}}
    :=
    a\left(
        \bigl(V^{q,r,\mathrm{in}}\bigr)_t^*
    \right).
\]
By \Cref{lem:adjusted-max},
$E^{q,r,\mathrm{in}}$ is a nondecreasing e-process valid under
every $R\in K_{q,r}^{\mathrm{in}}$.

We briefly verify preservation of the simultaneous growth rate.
Fix
\[
    Q\notin K_{q,r}^{\mathrm{in}}.
\]
Since $K_{q,r}^{\mathrm{in}}$ is weakly compact and
$Q\notin K_{q,r}^{\mathrm{in}}$,
\[
    \Phi_{K_{q,r}^{\mathrm{in}}}(Q)>0.
\]
Furthermore,
\[
    \bigl(V^{q,r,\mathrm{in}}\bigr)_t^*
    \ge
    V_t^{q,r,\mathrm{in}},
\]
so
\[
    \liminf_{t\to\infty}
    \frac1t
    \log
    \bigl(V^{q,r,\mathrm{in}}\bigr)_t^*
    \ge
    \Phi_{K_{q,r}^{\mathrm{in}}}(Q)
    >0,
    \qquad
    Q^\infty\text{-a.s.}
\]
In particular,
\[
    \log
    \bigl(V^{q,r,\mathrm{in}}\bigr)_t^*
    \longrightarrow\infty
\]
almost surely under $Q^\infty$. Since $a$ is growth preserving,
\[
    \frac{\log a(e^y)}{y}\longrightarrow1
    \qquad\text{as }y\to\infty.
\]
It follows that
\[
    \liminf_{t\to\infty}
    \frac1t
    \log E_t^{q,r,\mathrm{in}}
    \ge
    \Phi_{K_{q,r}^{\mathrm{in}}}(Q),
    \qquad
    Q^\infty\text{-a.s.}
\]
The process $E^{q,r,\mathrm{in}}$ is independent of $Q$.
Therefore $K_{q,r}^{\mathrm{in}}$ is simultaneously
REGROW-regular over
\[
    \mathcal P\setminus K_{q,r}^{\mathrm{in}}.
\]
This proves Condition~\textup{(ii)}.

\medskip

\noindent
\emph{Condition \textup{(iii)}: simultaneous exterior
witnesses.}

Fix $q,r\in\mathbb Q$ with $q<r$, and write
\[
    K_{q,r}^{\mathrm{out}}
    :=
    B_{q,r}^c
    =
    \{P_\theta:\theta\le q\}
    \cup
    \{P_\theta:\theta\ge r\},
\]
where the complement is taken relative to $\mathcal P$.

We construct a single e-process, depending only on $(q,r)$,
which is valid under every distribution in
$K_{q,r}^{\mathrm{out}}$ and achieves the REGROW rate under
every $P_{\theta_0}\in B_{q,r}$.

Enumerate the countable dense subset
\[
    \mathbb Q\cap(q,r)
    =
    \{h_m:m\ge1\},
\]
and set
\[
    \rho_m:=2^{-m},
    \qquad m\ge1.
\]
Then
\[
    \rho_m>0,
    \qquad
    \sum_{m=1}^\infty\rho_m=1.
\]

For each $h\in(q,r)$, define
\begin{align}
    L_t^{(q,h)}
    &:=
    \prod_{i=1}^t
    \frac{p_h(X_i)}{p_q(X_i)}
    =
    \exp\left\{
        (h-q)\sum_{i=1}^tX_i
        -
        \frac t2(h^2-q^2)
    \right\},
    \label{eq:gaussian-left-lr}\\
    L_t^{(r,h)}
    &:=
    \prod_{i=1}^t
    \frac{p_h(X_i)}{p_r(X_i)}
    =
    \exp\left\{
        (h-r)\sum_{i=1}^tX_i
        -
        \frac t2(h^2-r^2)
    \right\},
    \label{eq:gaussian-right-lr}
\end{align}
with
\[
    L_0^{(q,h)}=L_0^{(r,h)}=1.
\]
Define
\[
    E_t^{(h)}
    :=
    \min\left\{
        L_t^{(q,h)},
        L_t^{(r,h)}
    \right\},
    \qquad
    E_0^{(h)}=1.
\]

We first verify that $E^{(h)}$ is an e-process for the
composite null $K_{q,r}^{\mathrm{out}}$.

If $X\sim P_\theta$, then
\begin{align*}
    \mathbb E_\theta
    \left[
        \frac{p_h(X)}{p_q(X)}
    \right]
    &=
    \exp\left\{
        (h-q)\theta
        +
        \frac12(h-q)^2
        -
        \frac12(h^2-q^2)
    \right\}\\
    &=
    \exp\{(h-q)(\theta-q)\}.
\end{align*}
Because $h>q$, this quantity is at most one whenever
$\theta\le q$. Hence
\[
    \mathbb E_\theta
    \left[
        L_t^{(q,h)}
        \mid
        \mathcal F_{t-1}
    \right]
    \le
    L_{t-1}^{(q,h)}
\]
for every $\theta\le q$. Thus $L^{(q,h)}$ is a nonnegative test
supermartingale under every $P_\theta$ with $\theta\le q$.

Similarly,
\[
    \mathbb E_\theta
    \left[
        \frac{p_h(X)}{p_r(X)}
    \right]
    =
    \exp\{(h-r)(\theta-r)\}.
\]
Because $h-r<0$, this quantity is at most one whenever
$\theta\ge r$. Therefore $L^{(r,h)}$ is a nonnegative test
supermartingale under every $P_\theta$ with $\theta\ge r$.

Let $\tau$ be an arbitrary stopping time. By optional stopping
for nonnegative supermartingales, with the usual Fatou argument
when $\tau$ is unbounded, if $\theta\le q$, then
\[
\begin{aligned}
    \mathbb E_\theta
    \left[
        E_\tau^{(h)}
    \right]
    &\le
    \mathbb E_\theta
    \left[
        L_\tau^{(q,h)}
    \right]\\
    &\le1.
\end{aligned}
\]
If $\theta\ge r$, then
\[
\begin{aligned}
    \mathbb E_\theta
    \left[
        E_\tau^{(h)}
    \right]
    &\le
    \mathbb E_\theta
    \left[
        L_\tau^{(r,h)}
    \right]\\
    &\le1.
\end{aligned}
\]
Thus, for every $h\in(q,r)$,
\[
    \sup_{P_\theta\in K_{q,r}^{\mathrm{out}}}
    \mathbb E_\theta
    \left[
        E_\tau^{(h)}
    \right]
    \le1
\]
for every stopping time $\tau$. Hence $E^{(h)}$ is an e-process
for the composite null $K_{q,r}^{\mathrm{out}}$.

Now define the fixed mixture
\begin{equation}
\label{eq:gaussian-universal-mixture}
    V_t^{q,r,\mathrm{out}}
    :=
    \sum_{m=1}^\infty
    \rho_m E_t^{(h_m)}.
\end{equation}
This process depends only on the interval $(q,r)$ and not on the
true interior parameter.

For every fixed $t$ and every realized
$(X_1,\ldots,X_t)$, the functions
\[
    h\longmapsto L_t^{(q,h)}
    \quad\text{and}\quad
    h\longmapsto L_t^{(r,h)}
\]
are continuous and bounded on the compact interval $[q,r]$.
Consequently, the series in
\eqref{eq:gaussian-universal-mixture} is finite.

Moreover, for every stopping time $\tau$ and every
$P_\theta\in K_{q,r}^{\mathrm{out}}$, Tonelli's theorem gives
\begin{align*}
    \mathbb E_\theta
    \left[
        V_\tau^{q,r,\mathrm{out}}
    \right]
    &=
    \sum_{m=1}^\infty
    \rho_m
    \mathbb E_\theta
    \left[
        E_\tau^{(h_m)}
    \right]\\
    &\le
    \sum_{m=1}^\infty\rho_m\\
    &=1.
\end{align*}
Thus $V^{q,r,\mathrm{out}}$ is an e-process under every
distribution in $K_{q,r}^{\mathrm{out}}$.

We next prove that this single mixture achieves the desired
growth rate under every $P_{\theta_0}\in B_{q,r}$. Fix an
arbitrary
\[
    \theta_0\in(q,r).
\]
For every fixed $h\in(q,r)$, the strong law of large numbers
gives
\[
    \frac1t\sum_{i=1}^tX_i
    \longrightarrow
    \theta_0,
    \qquad
    P_{\theta_0}^\infty\text{-a.s.}
\]
It follows from \eqref{eq:gaussian-left-lr} that
\begin{align*}
    \frac1t\log L_t^{(q,h)}
    &\longrightarrow
    (h-q)\theta_0-\frac12(h^2-q^2)\\
    &=
    \frac12(\theta_0-q)^2
    -
    \frac12(\theta_0-h)^2,
\end{align*}
and from \eqref{eq:gaussian-right-lr} that
\begin{align*}
    \frac1t\log L_t^{(r,h)}
    &\longrightarrow
    (h-r)\theta_0-\frac12(h^2-r^2)\\
    &=
    \frac12(\theta_0-r)^2
    -
    \frac12(\theta_0-h)^2,
\end{align*}
almost surely under $P_{\theta_0}^\infty$. Therefore,
\begin{equation}
\label{eq:gaussian-component-rate}
    \frac1t\log E_t^{(h)}
    \longrightarrow
    I_{q,r}(\theta_0)
    -
    \frac12(\theta_0-h)^2,
    \qquad
    P_{\theta_0}^\infty\text{-a.s.},
\end{equation}
where
\begin{align}
    I_{q,r}(\theta_0)
    &:=
    \min\left\{
        \frac12(\theta_0-q)^2,
        \frac12(\theta_0-r)^2
    \right\}
    \nonumber\\
    &=
    \inf_{P_\theta\in K_{q,r}^{\mathrm{out}}}
    D_{\mathrm{KL}}(P_{\theta_0}\|P_\theta)
    \nonumber\\
    &=
    \Phi_{K_{q,r}^{\mathrm{out}}}(P_{\theta_0}).
    \label{eq:gaussian-exterior-kl}
\end{align}
Since the set $\{h_m:m\ge1\}$ is countable,
\eqref{eq:gaussian-component-rate} holds simultaneously for all
$m$ on an event of $P_{\theta_0}^\infty$-probability one.

For each fixed $m$,
\[
    V_t^{q,r,\mathrm{out}}
    \ge
    \rho_m E_t^{(h_m)}.
\]
Hence
\begin{align*}
    \liminf_{t\to\infty}
    \frac1t
    \log V_t^{q,r,\mathrm{out}}
    &\ge
    \lim_{t\to\infty}
    \left\{
        \frac{\log\rho_m}{t}
        +
        \frac1t\log E_t^{(h_m)}
    \right\}\\
    &=
    I_{q,r}(\theta_0)
    -
    \frac12(\theta_0-h_m)^2
\end{align*}
almost surely under $P_{\theta_0}^\infty$.

Because this inequality holds for every $m$,
\[
\begin{aligned}
    \liminf_{t\to\infty}
    \frac1t
    \log V_t^{q,r,\mathrm{out}}
    &\ge
    \sup_{m\ge1}
    \left\{
        I_{q,r}(\theta_0)
        -
        \frac12(\theta_0-h_m)^2
    \right\}\\
    &=
    I_{q,r}(\theta_0),
\end{aligned}
\]
where the last equality follows from density of
$\{h_m:m\ge1\}$ in $(q,r)$. Together with
\eqref{eq:gaussian-exterior-kl}, this gives
\[
    \liminf_{t\to\infty}
    \frac1t
    \log V_t^{q,r,\mathrm{out}}
    \ge
    \Phi_{K_{q,r}^{\mathrm{out}}}(P_{\theta_0}),
    \qquad
    P_{\theta_0}^\infty\text{-a.s.}
\]

The process $V^{q,r,\mathrm{out}}$ is fixed independently of
$\theta_0$, and the preceding conclusion holds for every
$\theta_0\in(q,r)$. Thus it is a simultaneous exterior witness,
except that it need not be nondecreasing.

To obtain a nondecreasing process, define
\[
    \bigl(V^{q,r,\mathrm{out}}\bigr)_t^*
    :=
    \max_{0\le u\le t}
    V_u^{q,r,\mathrm{out}},
\]
and
\[
    E_t^{q,r,\mathrm{out}}
    :=
    a\left(
        \bigl(V^{q,r,\mathrm{out}}\bigr)_t^*
    \right),
\]
where $a$ is the same fixed growth-preserving adjuster as above.
By \Cref{lem:adjusted-max},
$E^{q,r,\mathrm{out}}$ is a nondecreasing e-process under every
$P_\theta\in K_{q,r}^{\mathrm{out}}$.

Finally, fix $\theta_0\in(q,r)$. Since
\[
    \Phi_{K_{q,r}^{\mathrm{out}}}(P_{\theta_0})
    =
    I_{q,r}(\theta_0)
    >0,
\]
we have
\[
    \liminf_{t\to\infty}
    \frac1t
    \log
    \bigl(V^{q,r,\mathrm{out}}\bigr)_t^*
    \ge
    I_{q,r}(\theta_0)>0.
\]
Therefore,
\[
    \log
    \bigl(V^{q,r,\mathrm{out}}\bigr)_t^*
    \longrightarrow\infty
\]
almost surely under $P_{\theta_0}^\infty$. The
growth-preserving property of $a$ then yields
\[
    \liminf_{t\to\infty}
    \frac1t
    \log E_t^{q,r,\mathrm{out}}
    \ge
    I_{q,r}(\theta_0)
    =
    \Phi_{K_{q,r}^{\mathrm{out}}}(P_{\theta_0}),
    \qquad
    P_{\theta_0}^\infty\text{-a.s.}
\]
Since the same process $E^{q,r,\mathrm{out}}$ works for every
$\theta_0\in(q,r)$,
$K_{q,r}^{\mathrm{out}}=B_{q,r}^c$ is simultaneously
REGROW-regular over $B_{q,r}$. This proves
Condition~\textup{(iii)}.

We have verified all three conditions of
\Cref{def:witness-basis}. Therefore,
\[
    \mathscr B
    =
    \left\{
        B_{q,r}:q,r\in\mathbb Q,\ q<r
    \right\}
\]
is a countable local REGROW witness basis for
$\mathcal P$.
\end{proof}

\end{document}
